\subsection{NEIGHBOURHOODS OF A POINT IN A TOPOLOGICAL GROUP}

Let \(\mathfrak{B}\) be the neighbourhood filter of the identity element \(e\) in a topological group G, and let \(a\) be any point of G. Since \(x \to ax\) and \(x \to xa\) are homeomorphisms, it follows that the neighbourhood filter of \(a\) is the family \(a\) . \(\mathfrak{B}\) of sets \(a\) . V, where V runs through \(\mathfrak{B}\) , and is also the family \(\mathfrak{B}\) . \(a\) of sets V. \(a\) . Thus we know the neighbourhood filter of any point of a topological group as soon as we know the neighbourhood filter of the identity element \(e\) of the group.

If we say that \(xy\) and \(x^{-1}\) are continuous at \(x = y = e\) , we obtain (Chapter I, § 2, no. 1):

\((\mathbf{GV}_{\mathbf{I}})\) . Given any \(\mathbf{U} \in \mathfrak{B}\) , there exists \(\mathbf{V} \in \mathfrak{B}\) such that \(\mathbf{V} \cdot \mathbf{V} \subset \mathbf{U}\) .

\((\mathbf{GV}_{\mathbf{II}})\) . Given any \(\mathbf{U}\in \mathfrak{B},\) we have \(\mathbf{U}^{-1}\in \mathfrak{B}\)

Every filter \(\mathfrak{B}\) on \(G\) which satisfies \((\mathrm{GV}_{\mathbb{I}})\) and \((\mathrm{GV}_{\mathbb{II}})\) also satisfies \((\mathrm{GV}_{\alpha})\) . Given any \(U \in \mathfrak{B}\) , there exists \(V \in \mathfrak{B}\) such that \(V \cdot V^{-1} \subset U\) . For by \((\mathrm{GV}_{\mathbb{I}})\) , there exists \(W \in \mathfrak{B}\) such that \(W \cdot W \subset U\) , and by \((\mathrm{GV}_{\mathbb{II}})\) there exists \(V \in \mathfrak{B}\) such that \(V \subset W \cap W^{-1}\) ; hence \(V^{-1} \subset W\) and therefore \(V \cdot V^{-1} \subset W \cdot W \subset U\) .

Conversely, if a filter \(\mathfrak{B}\) on \(\mathbf{G}\) satisfies \((\mathrm{GV}_{\alpha})\) , it follows first of all that \(e\) belongs to every set \(\mathbf{U} \in \mathfrak{B}\) ; for if \(\mathbf{V} \in \mathfrak{B}\) is such that \(\mathbf{V} \cdot \mathbf{V}^{-1} \subset \mathbf{U}\) , then, since V is not empty, we have \(x.x^{-1}=e\in U\) for every \(x\in V\) . The condition \((\mathrm{GV}_{\alpha})\) therefore implies that \(V^{-1}\subset V.V^{-1}\subset U\) , so that \(U^{-1}\in\mathfrak{B}\) whenever \(U\in\mathfrak{B}\) . Finally, if \(V\in\mathfrak{B}\) is such that \(V.V^{-1}\subset U\) , and \(W\in\mathfrak{B}\) is such that \(W\subset V\cap V^{-1}\) , we have \(W.W\subset U\) . We see thus that \((\mathrm{GV}_{\alpha})\) is equivalent to the conjunction of \((\mathrm{GV}_{I})\) and \((\mathrm{GV}_{II})\) .

Finally, since \(x \rightarrow axa^{-1}\) is a homeomorphism which leaves e fixed, B has the following property:

\((\mathrm{GV}_{\mathrm{III}})\) . For all \(a \in \mathbf{G}\) and all \(\mathbf{V} \in \mathfrak{B}\) , we have \(a. \mathbf{V}.a^{-1} \in \mathfrak{B}\) .

These three properties of the filter \(\mathfrak{B}\) are characteristic:

PROPOSITION 1. Let G be a group and let B be a filter on G satisfying the axioms \((\mathrm{GV}_{\mathrm{I}})\) , \((\mathrm{GV}_{\mathrm{II}})\) and \((\mathrm{GV}_{\mathrm{III}})\) . Then there is a unique topology on G, compatible with the group structure of G, for which B is the neighbourhood filter of the identity element e. For this topology the neighbourhood filter of any point \(a \in G\) is the same as each of the two filters a.B and B.a.

If there is a topology with the required properties, then by what has been said above the neighbourhood filter of \(a\) coincides with each of the filters \(a\) . \(\mathfrak{B}\) and \(\mathfrak{B}.a\) ; hence the topology is unique, if it exists. Its existence will be established if we show 1) that the filters \(a\) . \(\mathfrak{B}\) are the neighbourhood filters of a topology on \(\mathbf{G}\) , and 2) that this topology is compatible with the group structure of \(\mathbf{G}\) .

\begin{enumerate}
\def\labelenumi{\arabic{enumi})}
\item
  The filter \(a. \mathfrak{B}\) satisfies axiom \((\mathrm{V}_{\mathrm{III}})\) (see Chapter I § 1, no. 2) by reason of \((\mathrm{GV}_{\mathrm{I}})\) and \((\mathrm{GV}_{\mathrm{II}})\) , as we have already seen; hence to show that \(a. \mathfrak{B}\) is the neighbourhood filter of \(a\) in a topology on \(G\) , we have to verify axiom \((\mathrm{V}_{\mathrm{IV}})\) . Let then \(V\) be any set of \(\mathfrak{B}\) , and \(W\) a set of \(\mathfrak{B}\) such that \(W.W \subset V\) ; then for any \(x \in a.W\) we have \(x.W \subset a.W.W \subset a.V\) , so that \(a.V\) belongs to the filter \(x. \mathfrak{B}\) ; hence \((\mathrm{V}_{\mathrm{IV}})\) is satisfied.
\item
  Let us now show that the topology defined by the neighbourhood filters \(a\) . \(\mathfrak{B}\) satisfies (GT'). Let \(a, b\) be any two points of G; if we put \(x = au\) and \(y = bv\) , then we have to show that \(xy^{-1}\) is as near as we please to \(ab^{-1}\) whenever \(u\) and \(v\) are close enough to \(e\) . Now \((ab^{-1})^{-1}(xy^{-1}) = buv^{-1}b^{-1}\) ; let U be any neighbourhood of \(e\) , then we shall have \(buv^{-1}b^{-1} \in U\) if \(uv^{-1} \in b^{-1}Ub = V\) , and \(V \in \mathfrak{B}\) by reason of \((\mathrm{GV}_{\mathrm{III}})\) . But by \((\mathrm{GV}_{\mathrm{I}})\) and \((\mathrm{GV}_{\mathrm{II}})\) there exists \(W \in \mathfrak{B}\) such that \(W.W^{-1} \subset V\) ; hence it is enough to take \(u \in W\) and \(v \in W\) in order to have \(xy^{-1} \in (ab^{-1})U\) . This completes the proof.
\end{enumerate}

A common method of defining a topology compatible with a group structure on G consists in giving a filter satisfying the axioms \((\mathrm{GV}_{\mathrm{I}})\) , \((\mathrm{GV}_{\mathrm{II}})\) and \((\mathrm{GV}_{\mathrm{III}})\) . The corresponding conditions for a filter base B are as follows:

\((\mathbf{GV}_{\mathbf{I}}^{\prime})\) . Given any \(\mathbf{U}\in \mathfrak{B}\) , there exists \(\mathbf{V}\in \mathfrak{B}\) such that \(\mathbf{V}.\mathbf{V}\subset \mathbf{U}\) . \((\mathrm{GV}_{\mathrm{II}}^{\prime})\) . Given any \(\mathbf{U} \in \mathfrak{B}\) , there exists \(\mathbf{V} \in \mathfrak{B}\) such that \(\mathbf{V}^{-1} \subset \mathbf{U}\) . \((\mathrm{GV}_{\mathrm{III}}^{\prime})\) . Given any \(a \in \mathbf{G}\) and any \(\mathbf{U} \in \mathfrak{B}\) , there exists \(\mathbf{V} \in \mathfrak{B}\) such that \(\mathbf{V} \subset a. \mathbf{U}. a^{-1}\) .

A neighbourhood of e which coincides with its image under the symmetry \(x \rightarrow x^{-1}\) is said to be symmetric. If V is any neighbourhood of e, then \(V \cup V^{-1}\) , \(V \cap V^{-1}\) and \(V \cdot V^{-1}\) are symmetric neighbourhoods. By \((\mathrm{GV}_{\mathrm{II}})\) , the symmetric neighbourhoods form a fundamental system of neighbourhoods of e. Also it follows from \((\mathrm{GV}_{\mathrm{I}})\) that when V runs through a fundamental system of neighbourhoods of e, the sets \(V^{n}\) (where n is a fixed integer \(\neq 0\) ) form a fundamental system of neighbourhoods of e.

Remark. If G is commutative, we have \(x \cdot A \cdot x^{-1} = A\) for every subset A of G and every \(x \in G\) , and therefore \((GV_{III})\) {[}resp. \((GV_{III}')\) {]} is automatically satisfied for every filter (resp. filter base) on G. On the other hand, if G is not abelian, then \((GV_{III})\) is not a consequence of \((GV_I)\) and \((GV_{II})\) {[}see Exercise 5{]}.

If G is a commutative group, written additively, the axioms which characterize the filter B of neighbourhoods of the origin for a topology compatible with the group structure of G are therefore the following: \((\mathrm{GA}_{\mathrm{I}})\) . Given any \(U \in B\) , there exists \(V \in B\) such that \(V + V \subset U\) . \((\mathrm{GA}_{\mathrm{II}})\) . Given any \(U \in B\) , we have \(-U \in B\) .

PROPOSITION 2. A topological group G is Hausdorff if and only if the set \(\{e\}\) is closed.

Clearly, if G is Hausdorff, then \(\{e\}\) is closed. Conversely, if \(\{e\}\) is closed, then the diagonal \(\Delta\) of \(G \times G\) is closed, because it is the inverse image of \(\{e\}\) under the continuous mapping \((x, y) \to xy^{-1}\) ; hence (Chapter I, § 8, no. 1, Proposition 1) G is Hausdorff.

COROLLARY. A topological group G is Hausdorff if and only if the intersection of the neighbourhoods of e consists only of the point e.

The condition is clearly necessary. Conversely, if the intersection of all the neighbourhoods of \(e\) is just \(\{e\}\) , then given any \(x \neq e\) there is a neighbourhood \(V\) of \(e\) such that \(x^{-1} \notin V\) and therefore \(e \notin xV\) . This shows that \(x\) is not in the closure of \(\{e\}\) , and thus \(\{e\}\) is closed, so that \(G\) is Hausdorff.

Example. Definition of a topology on a group by means of a set of subgroups. If \(\mathfrak{B}\) is a filter base on a group \(\mathbf{G}\) , formed of subgroups of \(\mathbf{G}\) , then it is immediately seen that \(\mathfrak{B}\) satisfies axioms \((\mathrm{GV}_{\mathrm{I}}^{\prime})\) and \((\mathrm{GV}_{\mathrm{II}}^{\prime})\) , since \(\mathbf{H}.\mathbf{H}^{-1} = \mathbf{H}\) for any subgroup \(\mathbf{H}\) of \(\mathbf{G}\) . Hence the set \(\mathfrak{B}\) will be a fundamental system of neighbourhoods of \(e\) in a topology compatible with the group structure of \(\mathbf{G}\) , provided that \(\mathfrak{B}\) satisfies \((\mathrm{GV}_{\mathrm{III}})\) ; this will in particular be the case if all the subgroups in \(\mathfrak{B}\) are normal, hence always if \(G\) is commutative. The topology thus defined is Hausdorff, by Proposition 2, if and only if the intersection of all the subgroups in \(\mathfrak{B}\) consists only of \(e\) . The most interesting cases are those in which the subgroup \(\{e\}\) is not in \(\mathfrak{B}\) (otherwise the topology defined by \(\mathfrak{B}\) is the discrete topology): if \(\{e\} \notin \mathfrak{B}\) , the topology defined by \(\mathfrak{B}\) is Hausdorff only if \(\mathfrak{B}\) is an infinite set.

Since the intersection of two subgroups is a subgroup, we can define a topology on G, compatible with its group structure, starting from any set \(\mathfrak{F}\) of subgroups of G: let \(\mathfrak{G}\) be the set of all subgroups \(a.H.a^{-1}\) , where \(H\in\mathfrak{F}\) and \(a\in G\) , and let \(\mathfrak{B}\) be the set of all finite intersections of subgroups belonging to \(\mathfrak{G}\) . Then \(\mathfrak{B}\) is a filter base and satisfies \((\mathrm{GV}_{\mathrm{III}}^{\prime})\) .

Consider in particular the additive group of a ring A. Every set \(\mathfrak{F}\) of ideals of A defines a topology compatible with this additive group structure. This topology is Hausdorff if the intersection of all the ideals of \(\mathfrak{F}\) is the zero ideal, and it is not discrete if no finite intersection of the ideals of \(\mathfrak{F}\) is the zero ideal. Topologies defined in this way play a large part in the theory of numbers (see the Exercises of §§ 6 and 7 of this chapter).

\subsection{ISOMORPHISMS AND LOCAL ISOMORPHISMS}

In accordance with the general definitions (Set Theory, Chapter IV, § 1, no. 5) an isomorphism \(f\) of a topological group \(G\) onto a topological group \(G'\) is a bijective mapping of \(G\) onto \(G'\) which is simultaneously an isomorphism of the group structure of \(G\) onto that of \(G'\) , and a homeomorphism of \(G\) onto \(G'\) . In other words, \(f\) is an isomorphism of \(G\) onto \(G'\) if and only if: 1) \(f\) is bijective; 2) \(f(xy) = f(x)f(y)\) for all \(x, y \in G\) ; and 3) \(f\) is bicontinuous.

For example, if \(a\) is any point of \(G\) , the mapping \(x \to axa^{-1}\) is an isomorphism of \(G\) onto \(G\) , that is (loc. cit.), an automorphism of the topological group \(G\) . It is called an inner automorphism.

If a topology \(\mathfrak{C}\) is compatible with the group structure of a group \(\mathbf{G}\) , then \(\mathfrak{C}\) is also compatible with the group structure opposite to that of \(\mathbf{G}\) . If \(\mathbf{G}^0\) denotes the topological group obtained by giving the group opposite to \(\mathbf{G}\) the topology \(\mathfrak{C}\) , then the symmetry \(x \to x^{-1}\) is an isomorphism of the topological group \(\mathbf{G}\) onto the topological group \(\mathbf{G}^0\) .

DEFINITION 2. If G and G' are two topological groups, a local isomorphism of G with G' is a homeomorphism f of a neighbourhood V of the identity element of G onto a neighbourhood V' of the identity element of G' which satisfies the following conditions:

\begin{enumerate}
\def\labelenumi{\arabic{enumi})}
\tightlist
\item
  For each pair \(x, y\) of points of \(V\) such that \(xy \in V\) ,
\end{enumerate}

\[
f (x y) = f (x) f (y).
\]

\begin{enumerate}
\def\labelenumi{\arabic{enumi})}
\setcounter{enumi}{1}
\tightlist
\item
  If \(g\) is the mapping inverse to \(f\) , then for each pair of points \(x', y'\) of \(V'\) such that \(x'y' \in V'\) , we have \(g(x'y') = g(x')g(y')\) .
\end{enumerate}

The mapping g is then a local isomorphism of \(G'\) with G.

Two topological groups G, G' are said to be locally isomorphic if there exists a local isomorphism of G with G'.

Isomorphic topological groups are evidently locally isomorphic. The converse is false.

* For example, we shall see in Chapter V, § 1, that the topological groups R and T are locally isomorphic but not isomorphic. *

If \(f\) is a local isomorphism of \(G\) with \(G'\) , then every restriction of \(f\) to a neighbourhood of the identity element of \(G\) is again a local isomorphism of \(G\) with \(G'\) .

A local isomorphism of G with G is called a local automorphism of G.

In general, if \(f\) is a homeomorphism of a neighbourhood \(V\) of the identity element of \(G\) onto a neighbourhood \(V'\) of the identity element of \(G'\) which satisfies condition 1) of Definition 2, \(f\) does not necessarily satisfy condition 2) (see Exercise 7). Nevertheless, \(G\) and \(G'\) are in fact locally isomorphic.

PROPOSITION 3. Let G and G' be two topological groups, and let f be a homeomorphism of a neighbourhood V of the identity element of G onto a neighbourhood V' of the identity element of G', which satisfies condition 1) of Definition 2. Then f is an extension of a local isomorphism of G with G'.

For it is not hard to see that if W is a neighbourhood of the identity element of G such that \(W.W \subset V\) , then the restriction of f to W is a local isomorphism of G with \(G'\) .

\section{SUBGROUPS, QUOTIENT GROUPS, HOMOMORPHISMS, HOMOGENEOUS SPACES, PRODUCT GROUPS}

\subsection{SUBGROUPS OF A TOPOLOGICAL GROUP}

Let G be a topological group and let H be a subgroup of G. By (GT'), the topology induced on H by the topology of G is compatible with the group structure of H. The structure of a topological group thus defined on H is said to be induced by that of G. Whenever we consider a subgroup H of G as a topological group, it is always this induced structure which is under consideration, unless the contrary is expressly stated.

PROPOSITION 1. The closure \(\overline{H}\) of a subgroup H of a topological group G is a subgroup of G. If H is a normal subgroup of G, then so is \(\overline{H}\) .

If \(a, b \in \overline{\mathbf{H}}\) then \(ab^{-1} \in \overline{\mathbf{H}}\) , because the mapping \((x, y) \to xy^{-1}\) is continuous on \(\mathbf{G} \times \mathbf{G}\) and transforms \(\mathbf{H} \times \mathbf{H}\) into \(\mathbf{H}\) (Chapter I, § 2, no. 1, Theorem 1). In the same way, the continuity of the mapping \(x \to axa^{-1}\) shows that if \(\mathbf{H}\) is normal then \(\overline{\mathbf{H}}\) is normal.

In particular, the closure N of the set \(\{e\}\) consisting only of the identity element of G, is a normal subgroup of G; and \(N=\{e\}\) if and only if G is Hausdorff (§ 1, no. 2, Proposition 2).

PROPOSITION 2. If G is a Hausdorff topological group, then the closure of a commutative subgroup of G is a commutative subgroup of G.

By reason of Proposition 1, we may limit ourselves to the case where H is dense in G. The continuous functions xy and yx are equal on H × H, and are therefore equal on G × G, by virtue of the principle of extension of identities (Chapter I, § 8, no. 1, Proposition 2, Corollary 1).

PROPOSITION 3. Let G be a Hausdorff topological group and let M be any subset of G. Then the set M' of elements of G which commute with each element of M is a closed subgroup of G. In particular, the centre of G is closed in G.

For \(\mathbf{M}'\) is the intersection of the sets \(\mathbf{F}_m(m \in \mathbf{M})\) , where \(\mathbf{F}_m\) is the set of all \(x \in \mathbf{G}\) such that \(xm = mx\) ; and \(\mathbf{F}_m\) is closed (Chapter I, § 8, no. 1, Proposition 2).

PROPOSITION 4. Let G be a topological group and let H be a subgroup of G which is locally closed at one point of H (Chapter I, § 3, no. 3, Definition 2). Then H is closed in G.

By translation, H is locally closed at each of its points, i.e.~H is locally closed in G. Let V be a symmetric open neighbourhood of e in G such that \(V \cap H\) is closed in V. If \(x \in \overline{H}\) , then xV meets H; if \(y \in xV \cap H\) , we have \(x \in yV\) , and \(y(V \cap H) = (yV) \cap H\) is closed in yV. But x lies in the closure of \((yV) \cap H\) , hence \(x \in H\) .

COROLLARY. A subgroup of a topological group is open if and only if it has an interior point. Every open subgroup is closed.

If a subgroup H has an interior point, then by translation all the points of H are interior and so H is open. The second assertion of the corollary is a particular case of Proposition 4.

PROPOSITION 5. A subgroup H of a topological group G is discrete if and only if H has an isolated point. Every discrete subgroup of a Hausdorff group is closed.

If H is discrete, every point of H is isolated. Conversely if H has an isolated point, then by translation every point of H is isolated and therefore H is discrete. If H is discrete and G is Hausdorff, then there is a neighbourhood V of e such that \(V \cap H = \{e\}\) ; \(\{e\}\) is closed in G and therefore a fortiori in V, so that H is locally closed at e. Hence H is closed in G by Proposition 4.

Remark. Let H be any subgroup of a topological group G. For each \(x \in \overline{H}\) we have \(\overline{xH} = x\) . \(\overline{H} = \overline{H}\) , since translation by x is a homeomorphism of G onto G. In other words, for each \(x \in \overline{H}\) , xH is dense in \(\overline{H}\) . It follows that if H is not closed, then \(\overline{H} \cap \mathrm{C}H\) is dense in \(\overline{H}\) .

\subsection{COMPONENTS OF A TOPOLOGICAL GROUP}

Let V be a symmetric neighbourhood of e in G. The subgroup generated by V, which is denoted by \(V^{\infty}\) consists of all products \(\prod_{i=1}^{n} x_{i}\) of finite sequences of elements of V. \(V^{\infty}\) is open, since it has e as an interior point, and is therefore closed by Proposition 4 of no. 1. It follows that:

PROPOSITION 6. A connected topological group is generated by every neighbourhood of the identity element.

The converse of this proposition is in general false, as we shall see in Chapter IV (§ 2, no. 5). If a topological group G is generated by each neighbourhood of the identity element, the most we can say is that G contains no open subgroup other than G.

* As an example of a non-connected group G which has an open subgroup distinct from G, we may cite the multiplicative group \(\mathbf{R}^{*}\) of non-zero real numbers, in which the subgroup \(\mathbf{R}_{+}^{*}\) of real numbers \textgreater0 is both open and closed (see Chapter IV, § 3, no. 2). *

PROPOSITION 7. In a topological group G, the component K of the identity element e is a closed normal subgroup. The component of any point \(x \in G\) is the coset \(x.K = K.x\) .

If \(a \in \mathbf{K}\) , then \(a^{-1}\mathbf{K}\) is connected and contains \(e\) ; hence \(\mathbf{K}^{-1}\mathbf{K} \subset \mathbf{K}\) , which shows that \(\mathbf{K}\) is a subgroup of \(\mathbf{G}\) . This subgroup is invariant under all automorphisms of \(\mathbf{G}\) , and in particular under all inner automorphisms, therefore \(\mathbf{K}\) is normal in \(\mathbf{G}\) ; also \(\mathbf{K}\) is closed (Chapter I, § 11, no. 5, Proposition 9). Finally, the left translation \(y \to xy\) is a homeomorphism of \(\mathbf{G}\) which sends \(e\) to \(x\) , and hence the component of \(x\) is \(x\cdot \mathbf{K}\) .

The component of the identity element e of G is called the identity component of G.

\subsection{DENSE SUBGROUPS}

The following Proposition generalizes Proposition 1 of no. 1:

PROPOSITION 8. Let H be a dense subgroup of a topological group G, and let K be a normal subgroup of H. Then the closure \(\overline{K}\) of K in G is a normal subgroup of G.

For the mapping \((z, x) \to zxz^{-1}\) is continuous on \(\mathbf{G} \times \mathbf{G}\) , and maps \(\mathbf{H} \times \overline{\mathbf{K}}\) into \(\overline{\mathbf{K}}\) ; hence (Chapter I, § 2, no. 1, Theorem 1) it maps \(\mathbf{G} \times \overline{\mathbf{K}} = \overline{\mathbf{H}} \times \overline{\mathbf{K}}\) into \(\overline{\mathbf{K}}\) .

PROPOSITION 9. Let H be a dense subgroup of a topological group G. If H is generated by every neighbourhood of the identity element in H, then G is generated by every neighbourhood of the identity element in G.

Let V be any symmetric neighbourhood of e in G. Then \(V \cap H\) is a neighbourhood of e in H, and hence generates H. It follows that V generates a subgroup \(H'\) which contains H; but \(H'\) is open and closed (no. 1, Corollary to Proposition 4), and therefore contains \(\overline{H} = G\) .

\subsection{SPACES WITH OPERATORS}

Let X be a topological space and let G be a topological group. G is said to operate continuously on X if the following conditions are satisfied:

\begin{enumerate}
\def\labelenumi{\arabic{enumi})}
\item
  X has G as a group of operators; in other words X is endowed with an external law of composition \((s, x) \to s.x\) for which G is the set of operators, and which is such that \(s.(t.x) = (st).x\) and e.x = x for all \(s, t \in G\) and all \(x \in X\) .
\item
  The mapping \((s, x) \to s.x\) of \(\mathbf{G} \times \mathbf{X}\) into \(\mathbf{X}\) is continuous.
\end{enumerate}

LEMMA 1. If a topological group G operates continuously on a topological space X, then for each \(s \in \mathbf{G}\) the mapping \(x \to s.x\) is a homeomorphism of X onto X.

For this mapping is a continuous bijection whose inverse \(x \rightarrow s^{-1}.x\) is also continuous.

We recall that for each \(x \in X\) the set G.x of transforms s.x of x by the elements s of G is called the orbit of x (with respect to the group of operators G), and that the set of all \(s \in G\) such that s.x = x is a subgroup of G called the stabilizer of x. The relation \(\mathrm{R}(x, y) : ``y \text{ belongs to the orbit of } x''\) is an equivalence relation on X, called the equivalence relation defined by G; the equivalence classes with respect to this relation are the orbits of the points of X. The topological space X/R is called the orbit space of X (with respect to G), or the quotient space of X by the group G, and is denoted by X/G; and the topology of X/G is said to be the quotient of the topology of X by G.

LEMMA 2. If a topological group G operates continuously on a topological space X, then the equivalence relation R defined by G is open.

For the saturation with respect to \(\mathbf{R}\) of an open subset \(\mathbf{U}\) of \(\mathbf{X}\) is the set \(\bigcup_{s\in G}s.\mathbf{U},\) and each \(s.\mathbf{U}\) is open by Lemma 1.

Examples. 1) Let H be a subgroup of a topological group G. H operates continuously on G by the external law \((s, x) \rightarrow sx\) . H also operates continuously on G by the external law \((x, s) \rightarrow sxs^{-1}\) .

\begin{enumerate}
\def\labelenumi{\arabic{enumi})}
\setcounter{enumi}{1}
\item
  *If K is a topological division ring (§ 6, no. 7), the multiplicative group K* operates continuously on K by the external law \((s, x) \to sx\) . *
\item
  Let G be a topological group, X a topological space. Then the mapping \((s, x) \to x\) of \(G \times X\) into X is an external law of composition on X, and G operates continuously on X with respect to this law; G is then said to operate trivially on X.
\end{enumerate}

Remark. Instead of saying that a topological group G operates continuously on a topological space X, it is often said that G operates continuously on the left on X. When the topological group \(G^{0}\) opposite to G operates continuously on X, we say that G operates continuously on the right on X. It comes to the same thing to say that X has a continuous external law of composition \((s, x) \to s.x\) with G as set of operators, such that \(s.(t.x) = (ts).x\) and e.x = x. Such a law is often written on the right: \((s, x) \to x.s\) (whence the terminology), and we have then \((x.t).s = x.(ts)\) . If G operates continuously on the right on X by the law \((s, x) \to x.s\) , then also G operates continuously on the left on X according to the external law \((s, x) \to x.s^{-1}\) , by virtue of axiom \((\mathrm{GT}_{\mathrm{II}})\) .

Let X (resp. X') be a set with a group of operators G (resp. G') and let \(f: \mathbf{G} \to \mathbf{G}'\) be a homomorphism and \(g: \mathbf{X} \to \mathbf{X}'\) a mapping. \(f\) and \(g\) are said to be compatible if \(g(s.x) = f(s).g(x)\) for all \(s \in \mathbf{G}\) and all \(x \in \mathbf{X}\) . If \(\mathbf{X}''\) is a third set with a group of operators G'', if \(f': \mathbf{G}' \to \mathbf{G}''\) is a homomorphism, \(g': \mathbf{X}' \to \mathbf{X}''\) a mapping, and if \(f'\) and \(g'\) are compatible, then \(f' \circ f\) and \(g' \circ g\) are compatible. When X, X' are topological spaces and G, G' are topological groups operating continuously on X, X' respectively, \((f, g)\) is said to be a morphism of the space with operators X into the space with operators X', provided that \(f\) and \(g\) are continuous and compatible. Passing to the quotients, \(g\) then induces a continuous mapping \(\mathbf{X}/\mathbf{G} \to \mathbf{X}'/\mathbf{G}'\) (Chapter I, § 3, no. 4, Corollary to Proposition 6).

Let G be a topological group operating continuously on a topological space X, and let \(\varphi\) be the canonical mapping of X onto the orbit space X/G. Let A be any subset of X, and let \(A'\) be the subspace of X which is the saturation of A with respect to the equivalence relation R defined by G (thus \(A'\) is the union of the orbits of points of A, and is said to be the saturation of A with respect to G). G operates continuously on \(A'\) by the restriction of \((s, x) \to s.x\) to \(G \times A'\) . Moreover, since R is open (Lemma 2) and \(A'\) is saturated, it follows from Proposition 4 of Chapter I, § 5, no. 2, and from the relation \(\varphi(A) = \varphi(A')\) , that:

PROPOSITION 10. The canonical bijection of the subspace \(\varphi(A)\) of \(X/G\) onto the orbit space \(A'/G\) is a homeomorphism.

Now let S be an equivalence relation on X such that, for each \(s \in G\) , the mapping \(x \to s.x\) is compatible with S {[}in other words, such that the relation \(x \equiv y \pmod{S}\) implies \(s.x \equiv s.y \pmod{S}\) {]}; for the sake of brevity we shall say that the relation S is compatible with the group G. If \(\psi\) is the canonical mapping of X onto X/S, and if \(s.\psi(x)\) denotes the class mod S of s.x, then X/S has G as groups of operators with respect to the external law \((s, \psi(x)) \to s.\psi(x) = \psi(s.x)\) . Moreover:

PROPOSITION 11. If the equivalence relation S on X is open and compatible with G, then G operates continuously on X/S.

Since the relation of equality on G and the relation S on X are open, it is enough to show that the mapping \((s,x)\to s.\psi(x)=\psi(s.x)\) of \(G\times X\) into X/S is continuous (Chapter I, § 5, no. 3, Corollary to Proposition 8); but this follows from the continuity of \(\psi\) and of the mapping

\[
(s, x) \rightarrow s. x.
\]

Remark. Let \(G'\) be another topological group operating continuously on X, and suppose that \(s.(s'x) = s'(s.x)\) for all \(s \in G\) , \(s' \in G'\) and \(x \in X\) . Then the equivalence relation \(S\) defined by \(G'\) is compatible with \(G\) , and since \(S\) is open (Lemma 2) it follows that \(G\) operates continuously on \(X/G'\) . Similarly, \(G'\) operates continuously on \(X/G\) . In these circumstances the two operations of the two groups \(G\) and \(G'\) on \(X\) are said to commute.

\subsection{HOMOGENEOUS SPACES}

Let G be a topological group and H a subgroup of G. H operates continuously on the right on G according to the external law \((t, x) \to xt\) , and the orbit of a point \(x \in G\) is the left coset xH. The set of orbits is therefore what we have called in algebra the homogeneous space G/H. Whenever we speak of G/H as a topological space, we shall always mean the orbit space of G (with respect to H) unless the contrary is expressly stated; i.e.~the quotient space of G by the equivalence relation \(x^{-1}y \in H\) . In conformity with the general definitions, we say that the topology of this space is the quotient by H of the topology of G.

PROPOSITION 12. The group G operates continuously on every homogeneous space G/H.

Since the equivalence relation \(x^{-1}y \in H\) is open (no. 4, Lemma 2) this is a particular case of Proposition 11 of no. 4.

PROPOSITION 13. Let G be a topological group and let H be a subgroup of G. Then the homogeneous space G/H is Hausdorff if and only if H is closed in G.

H is an equivalence class for the relation \(x^{-1}y \in H\) and therefore, if G/H is Hausdorff, H is closed in G. Conversely, if H is closed, then the graph of this relation is closed in \(G \times G\) , since it is the inverse image of H under the continuous mapping \((x, y) \to x^{-1}y\) . Since the relation \(x^{-1}y \in H\) is open, it follows from Chapter I, § 8, no. 3, Proposition 8 that G/H is Hausdorff.

PROPOSITION 14. Let G be a topological group and let H be a subgroup of G. Then the homogeneous space G/H is discrete if and only if H is open in G.

For the inverse images in G of the points of G/H under the canonical mapping are the cosets xH (x ∈ G); and these sets are open in G if and only if H is open in G.

Let X be a topological space on which a topological group G operates continuously and transitively; X is then (in the algebraic sense) a homogeneous space of G. Let \(x\) be a point of \(X\), \(H_x\) its stabilizer. The continuous surjection \(s \to s.x\) of \(G\) onto \(X\) factorizes canonically as

\[
\mathbf {G} \xrightarrow {f _ {x}} \mathbf {G} / \mathbf {H} _ {x} \xrightarrow {g _ {x}} \mathbf {X}
\]

\pandocbounded{\includegraphics[keepaspectratio]{images/part002_c9cfa224.jpg}}

where \(f_x\) is the canonical mapping of \(G\) onto the homogeneous space \(G/H_x\) , and \(g_x\) is the bijection \(s.H_x \to s.x\) of \(G/H_x\) onto \(X\) ; moreover (Chapter I, § 3, no. 4, Proposition 6) \(g_x\) is a continuous mapping. But \(g_x\) is not necessarily a homeomorphism of \(G/H_x\) onto \(X\) (Exercise 29). If \(g_x\) is a homeomorphism for each \(x \in X\) , then we say that \(X\) is a topological homogeneous space (of the topological group \(G\) ); for this to be so it is necessary and sufficient that, for each \(x \in E\) , the mapping \(s \to s.x\) of \(G\) into \(X\) should be open.

PROPOSITION 15. Let X be a topological space on which a topological group G operates continuously and transitively. For X to be a topological homogeneous space (relative to G) it is sufficient that for some point \(x_{0} \in X\) the mapping \(s \to s \cdot x_{0}\) transforms each neighbourhood of e in G into a neighbourhood of \(x_{0}\) in X.

Every \(x \in X\) can be written as \(x = t.x_0\) for some \(t \in G\) . If \(V\) is a neighbourhood of \(e\) , then \(V.x = (Vt).x_0\) is a neighbourhood of \(x\) , for we can write \((Vt).x_0 = t((t^{-1}Vt).x_0)\) , and the assertion follows from the facts that \(t^{-1}Vt\) is a neighbourhood of \(e\) in \(G\) and that \(y \to t.y\) is a homeomorphism of \(X\) onto itself (no. 4, Lemma 1). It follows that if \(U\) is any open subset of \(G\) and if \(x\) is any point of \(X\) , then \(U.x\) is open in \(X\) ; for if \(t \in U\) , then \(t^{-1}U\) is a neighbourhood of \(e\) , hence \((t^{-1}U).x\) is a neighbourhood of \(x\) , and \(t.((t^{-1}U).x) = U.x\) is a neighbourhood of \(t.x\) . Hence \(U.x\) is open in \(X\) , so that the mapping \(s \to s.x\) of \(G\) into \(X\) is open. This completes the proof.

\subsection{QUOTIENT GROUPS}

PROPOSITION 16. Let G be a topological group and let H be a normal subgroup of G. Then the quotient by H of the topology of G is compatible with the group structure of G/H.

If \(x \to \dot{x}\) is the canonical mapping of G onto G/H, then we have to show that \((\dot{x}, \dot{y}) \to \dot{x}\dot{y}^{-1}\) is a continuous mapping of (G/H) \(\times\) (G/H) into G/H. Since the equivalence relation \(x^{-1}y \in H\) is open (no. 4, Lemma 2), it is enough to show that \((x,y) \to \dot{x}\dot{y}^{-1}\) is a continuous mapping of G \(\times\) G into G/H (Chapter I, § 5, no. 3, Corollary to Proposition 8, and § 3, no. 4, Proposition 6). But \((x,y) \to \dot{x}\dot{y}^{-1}\) is the composition of the continuous mappings \(x \to \dot{x}\) and \((x,y) \to xy^{-1}\) , hence is continuous.

Whenever in the sequel we consider a quotient group G/H of a topological group G as a topological group, it is always to be understood that the topology of G/H is the quotient by H of the topology of G, unless the contrary is expressly stated.

PROPOSITION 17. Let \(\varphi\) be the canonical mapping of a topological group \(\mathbf{G}\) onto a quotient group \(\mathrm{G / H}\) . If \(\mathfrak{B}\) is a fundamental system of neighbourhoods of \(e\) in \(\mathbf{G}\) , then \(\varphi(\mathfrak{B})\) is a fundamental system of neighbourhoods of the identity element \(\varphi(e)\) of \(\mathrm{G / H}\) .

This is a particular case of Proposition 5 of Chapter I, § 5, no. 3.

Propositions 13 and 14 give in particular, for quotient groups:

PROPOSITION 18. Let G be a topological group and let H be a normal subgroup of G.

\begin{enumerate}
\def\labelenumi{\alph{enumi})}
\item
  The quotient group G/H is Hausdorff if and only if H is closed in G.
\item
  The quotient group G/H is discrete if and only if H is open in G.
\end{enumerate}

If G is a topological group and N is the closure of \(\{e\}\) in G, then N is a closed normal subgroup of G (no. 1, Proposition 1), hence G/N is Hausdorff; G/N is called the Hausdorff group associated with G.

PROPOSITION 19. If H is a discrete normal subgroup of a topological group G, then G/H is locally isomorphic to G.

Let \(\mathbf{V}\) be a neighbourhood of \(e\) in \(\mathbf{G}\) which does not contain any point of \(\mathbf{H}\) other than \(e\) , and let \(\mathbf{W}\) be a symmetric open neighbourhood of \(e\) in \(\mathbf{G}\) such that \(\mathbf{W}^2 \subset \mathbf{V}\) . Then the restriction to \(\mathbf{W}\) of the canonical mapping \(\varphi\) of \(\mathbf{G}\) onto \(\mathbf{G} / \mathbf{H}\) is injective; for if \(x, y \in \mathbf{W}\) and \(\varphi(x) = \varphi(y)\) , then \(x^{-1}y \in \mathbf{W}^2 \subset \mathbf{V}\) and \(x^{-1}y \in \mathbf{H}\) , so that \(x = y\) . By Proposition 17 it follows that the restriction of \(\varphi\) to \(\mathbf{W}\) is a homeomorphism of \(\mathbf{W}\) onto \(\varphi(\mathbf{W})\) ; since moreover \(\varphi(xy) = \varphi(x)\varphi(y)\) for all \(x, y \in \mathbf{W}\) , we conclude that \(\mathbf{G}\) and \(\mathbf{G} / \mathbf{H}\) are locally isomorphic (§ 1, no. 3, Proposition 3).

\subsection{SUBGROUPS AND QUOTIENT GROUPS OF A QUOTIENT GROUP}

Let G be a topological group, let H be a normal subgroup of G, and let \(\varphi: G \to G/H\) be the canonical mapping. We know that if \(A'\) is a subgroup of G/H, then \(\varphi^{-1}(A')\) is a subgroup of G which contains H. Conversely, if A is a subgroup of G, then \(\varphi(A)\) is a subgroup of G/H; moreover, there is a canonical bijection of the quotient group \(A/(A \cap H)\) onto the subgroup \(\varphi(A)\) of G/H, and a canonical bijection of \(\varphi(A)\) onto the quotient group AH/H, and both these bijections are isomorphisms for the group structures.

PROPOSITION 20. Let A be a subgroup of a topological group G, let H be a normal subgroup of G, and let \(\varphi\) denote the canonical mapping of G onto G/H. Then the canonical bijection of \(\varphi(A)\) onto AH/H is an isomorphism of topological groups.

This follows from the preceding remarks, and from Proposition 10 of no. 4.

\pandocbounded{\includegraphics[keepaspectratio]{images/part002_0097f9c2.jpg}}

The canonical bijection of \(\mathbf{A}/(\mathbf{A} \cap \mathbf{H})\) onto \(\varphi(\mathbf{A})\) is a continuous homomorphism, since it arises from the restriction of \(\varphi\) to \(\mathbf{A}\) by passing to the quotients; but in general the topological groups \(\mathbf{A}/(\mathbf{A} \cap \mathbf{H})\) and \(\mathbf{AH}/\mathbf{H}\) are not isomorphic (cf.~§ 4, no. 1, Proposition 1, Corollary 3).

* For example, take G to be the additive group R of real numbers, take H to be the group Z of integers, and take A to be the group θZ of integer multiples of an irrational number θ. Then A∩H = \{0\}, so that A/(A∩H) is a discrete group, isomorphic to Z; on the other hand, A + H is dense in R (as we shall see in Chapter V, § 1, no. 1, Proposition 1), hence (A + H)/H, which is locally isomorphic to A + H (no. 6, Proposition 19), is not a discrete group and consequently is not isomorphic to A/(A∩H). *

Nevertheless, we have the following proposition:

PROPOSITION 21. Let G be a topological group, \(G_{0}\) a dense subgroup of G, \(H_{0}\) a normal subgroup of \(G_{0}\) , H the closure of \(H_{0}\) in G and \(\varphi\) the canonical mapping \(G \rightarrow G/H\) . Then the canonical bijection \(G_{0}/H_{0} \rightarrow \varphi(G_{0})\) is an isomorphism of the topological group \(G_{0}/H_{0}\) onto a dense subgroup of G/H.

Since \(\mathbf{H}_0 = \mathbf{H} \cap \mathbf{G}_0\) , it is enough to show that if \(\mathbf{U}_0\) is any open subset of \(\mathbf{G}_0\) which is saturated with respect to the relation \(x^{-1}y \in \mathbf{H}_0\) , then \(\mathbf{U}_0\) is the intersection of \(\mathbf{G}_0\) and an open subset of \(\mathbf{G}\) which is saturated with respect to the relation \(x^{-1}y \in \mathbf{H}\) (Chapter I, § 3, no. 6, Proposition 10). Let \(\mathbf{U}\) be an open subset of \(\mathbf{G}\) such that \(\mathbf{U}_0 = \mathbf{U} \cap \mathbf{G}_0\) . Since \(\mathbf{U}_0 = \mathbf{U}_0\mathbf{H}_0\) , it is easily seen that \(\mathbf{U}_0 = \mathrm{UH}_0 \cap \mathbf{G}_0\) ; but \(\mathrm{UH}_0\) is open in \(\mathbf{G}\) , so we may suppose that \(\mathbf{U} = \mathrm{UH}_0\) . The set \(\mathrm{UH}\) is open in \(\mathbf{G}\) and is saturated with respect to the relation \(x^{-1}y \in \mathbf{H}\) ; hence the proposition will be proved if we show that \(\mathrm{UH} \cap \mathbf{G}_0 = \mathbf{U}_0\) . Now, if \(u \in \mathbf{U}\) and \(h \in \mathbf{H}\) are such that \(uh \in \mathbf{G}_0\) , then there is a symmetric neighbourhood \(\mathbf{V}\) of \(e\) in \(\mathbf{G}\) such that \(uv \subset U\) ; since \(Vh\) is a neighbourhood of \(h\) in \(\mathbf{G}\) , there exists \(z \in V\) such that \(zh \in \mathbf{H}_0\) . But then we have \(uz^{-1} \in U\) , and \(uh = (uz^{-1})(zh)\) ; therefore \(\mathrm{UH} \cap \mathbf{G}_0 \subset \mathrm{UH}_0\) . But \(\mathrm{UH}_0 = U\) , hence \(\mathrm{UH} \cap \mathbf{G}_0 \subset U \cap \mathbf{G}_0 = U_0\) . This completes the proof.

Let G be a topological group operating continuously on a topological space X, and let K be a normal subgroup of G which is contained in the stabilizer of each point of X. For each \(x \in X\) , the relation \(s \equiv t \pmod{K}\) therefore implies s.x = t.x, and passing to the quotient we define a mapping \(\dot{s} \to \dot{s}.x\) of G/K into X. It is immediately verified that with respect to the external law \((\dot{s}, x) \to \dot{s}.x\) , X has G/K as a group of operators. Moreover, G/K operates continuously on X with respect to this law; for the relation of equality on X and the relation \(s \equiv t \pmod{K}\) on G are both open equivalence relations, and therefore the result follows from the continuity of the mapping

\[
(s, x) \rightarrow \dot {s}. x = s. x
\]

of \(\mathbf{G} \times \mathbf{X}\) into \(\mathbf{X}\) (Chapter I, § 5, no. 3, Corollary to Proposition 8, and § 3, no. 4, Proposition 6).

Now let G be a topological group operating continuously on a topological space X, and let H be any normal subgroup of G; then H operates continuously on X. Let S be the equivalence relation defined by H on X; then S is open (no. 4, Lemma 2). The relation S is compatible with the group G (no. 4); for if \(y \equiv x \pmod{S}\) then there exists \(t \in H\) such that y = t.x; hence for all \(s \in G\) we have \(s.y = (sts^{-1}) \cdot s.x\) , and \(sts^{-1} \in H\) since H is normal in G; thus \(s.y \equiv s.x \pmod{S}\) . If \(\psi\) is the canonical mapping of X onto X/S, the group G therefore operates continuously on X/S with respect to the external law

\[
(s, \psi (x)) \rightarrow \psi (s. x)
\]

(no. 4, Proposition 11). Moreover, the group H is contained in the stabilizer of each of the points of X/S; as we have seen above, G/H operates continuously on X/S = X/H with respect to the external law \((s, \psi(x)) \to \psi(s.x)\) . If R denotes the equivalence relation on X defined by G, then the relation S implies R, and the equivalence relation R/S on X/S is that defined by the group G/H. Hence (Chapter I, § 3, no. 4, Proposition 7):

PROPOSITION 22. Let G be a topological group operating continuously on a topological space X, and let H be a normal subgroup of G. Then the canonical bijection of X/G onto (X/H)/(G/H) is a homeomorphism.

COROLLARY. Let G be a topological group, H a normal subgroup of G, and K a normal subgroup of G which contains H. Then the canonical bijection of G/K onto (G/H)/(K/H) is an isomorphism of topological groups.

We know already that this bijection is an isomorphism of groups, and Proposition 22 (applied to the group K operating on the right on G) shows that it is a homeomorphism.

\subsection{CONTINUOUS HOMOMORPHISMS AND STRICT MORPHISMS}

PROPOSITION 23. A homomorphism \(f\) of a topological group \(\mathbf{G}\) into a topological group \(\mathbf{G}'\) is continuous on \(\mathbf{G}\) if and only if it is continuous at one point of \(\mathbf{G}\) .

Suppose \(f\) is continuous at a point \(a \in \mathbf{G}\) ; then if \(\mathbf{V}'\) is any neighbourhood of \(f(a)\) , \(\mathbf{V} = f^{-1}(\mathbf{V}')\) is a neighbourhood of \(a\) . Hence if \(x\) is any point of \(\mathbf{G}\) , we have

\[
f (x a ^ {- 1} \mathrm{V}) = f (x) [ f (a) ] ^ {- 1} f (\mathrm{V}) \subset f (x) [ f (a) ] ^ {- 1} \mathrm{V} ^ {\prime},
\]

and therefore \(f\) is continuous at \(x\) .

A continuous homomorphism of a topological group G into a topological group \(G'\) is also called a morphism of G into \(G'\) for the topological group structures.

Let \(f\) be a continuous homomorphism of a topological group \(G\) into a topological group \(G'\) ; the inverse image \(H = f^{-1}(e')\) of the identity element \(e'\) of \(G'\) is a normal subgroup of \(G\) , and \(f(G)\) is a subgroup of \(G'\) . Consider the canonical factorization \(f = \psi \circ \overline{f} \circ \varphi\) , where \(\varphi\) is the canonical mapping \(G \to G/H\) , \(\psi\) is the canonical injection \(f(G) \to G'\) and lastly \(\overline{f}\) is a continuous bijective homomorphism of the quotient group \(G/H\) onto the subgroup \(f(G)\) (Chapter I, § 3, no. 5); \(\overline{f}\) is said to be the bijective homomorphism associated with \(f\) . In general, \(f\) is not an isomorphism of topological groups.

For example, let \(G'\) be a non-discrete topological group, and G the topological group obtained by giving \(G'\) the discrete topology; then the identity mapping of G into \(G'\) is a continuous bijective homomorphism, but is not bicontinuous.

DEFINITION 1. A continuous homomorphism of a topological group G into a topological group \(\mathbf{G}'\) is said to be a strict morphism of G into \(\mathbf{G}'\) if the bijective homomorphism \(\overline{f}\) of \(\mathbf{G}/f^{-1}(e')\) onto \(f(\mathbf{G})\) , associated with \(f\) , is an isomorphism of topological groups (in other words, if \(\overline{f}\) is bicontinuous).

An isomorphism of a topological group G onto a topological group \(G'\) is therefore a bijective strict morphism of G onto \(G'\) .

PROPOSITION 24. Let \(f\) be a continuous homomorphism of a topological group \(G\) into a topological group \(G'\) . Then the following three statements are equivalent: a) \(f\) is a strict morphism.

\begin{enumerate}
\def\labelenumi{\alph{enumi})}
\setcounter{enumi}{1}
\item
  The image under \(f\) of every open set in \(\mathbf{G}\) is an open set in \(f(\mathbf{G})\) .
\item
  The image under \(f\) of every neighbourhood of the identity element in \(G\) is a neighbourhood of the identity element in \(G'\) .
\end{enumerate}

In view of Lemma 2 of no. 4, the equivalence of a) and b) follows immediately from the definitions (Chapter I, § 5, no. 3, Proposition 5). The equivalence of b) and c) is a particular case of Proposition 15 of no. 5, if we observe that G operates continuously on \(f(\mathbf{G})\) by the external law \((s, f(t)) \to f(st)\) .

Remarks. 1) From condition b) of Proposition 24 it follows that every continuous homomorphism of a topological group into a discrete group is a strict morphism.

If G is compact and \(f(\mathbf{G})\) is Hausdorff, then the bijective homomorphism \(\overline{f}\) associated with f is bicontinuous (Chapter I, § 10, no. 2, Theorem 1, Corollary 2, and no. 1, Proposition 5, Corollary 4). Hence every continuous homomorphism of a compact group into a Hausdorff group is a strict morphism.

\pandocbounded{\includegraphics[keepaspectratio]{images/part002_802da4a7.jpg}}

\begin{enumerate}
\def\labelenumi{\arabic{enumi})}
\setcounter{enumi}{1}
\item
  Let \(f\) be a strict morphism of \(G\) into \(G'\) and let \(g\) be a strict morphism of \(G'\) into \(G''\) . If \(f\) is surjective or if \(g\) is injective, it follows immediately from Proposition 24 that \(g \circ f\) is a strict morphism of \(G\) into \(G''\) . But this conclusion is no longer necessarily valid if neither of the two preceding conditions is satisfied, even if \(f\) is injective and \(g\) surjective (Exercise 19).
\item
  Let \(f\) be a continuous homomorphism of a topological group \(G\) into a topological group \(G'\) , and let \(H\) be a normal subgroup of \(G\) . \(f\) induces a homomorphism \(g\) of the group \(G / H\) onto the quotient group \(f(G) / f(H)\) . This homomorphism \(g\) is continuous. Moreover, if \(f\) is a strict morphism of \(G\) into \(G'\) , then \(g\) is a strict morphism of \(G / H\) onto \(f(G) / f(H)\) ; for if \(U\) is open in \(G / H\) , and if \(\varphi\) (resp. \(\varphi'\) ) denotes the canonical mapping of \(G\) onto \(G / H\) {[}resp. of \(f(G)\) onto \(f(G) / f(H)\) {]}, then we have \(g(u) = \varphi'(f(\varphi^{-1}(u)))\) , and since \(\varphi^{-1}(u)\) is open in \(G\) , it follows that \(g(u)\) is open in \(f(G) / f(H)\) , which proves our assertion.
\end{enumerate}

\subsection{PRODUCTS OF TOPOLOGICAL GROUPS}

Let \((\mathbf{G}_i)_{i\in I}\) be a family of topological groups. Then we can define a group structure on the product set

\[
\mathbf {G} = \prod_ {i \in I} \mathbf {G} _ {i}
\]

(the product of the group structures of the \(\mathbf{G}_i\) ) by defining \((x_i) \cdot (y_i) = (x_i y_i)\) . If \(e_i\) is the identity element of \(\mathbf{G}_i\) , then \(e = (e_i)\) is the identity element of \(\mathbf{G}\) , and we have \((x_i)^{-1} = (x_i^{-1})\) . The product topology (Chapter I, § 2, no. 3) of the topologies of the \(\mathbf{G}_i\) is compatible with this group structure. For the mapping \(((x_i), (y_i)) \to (x_i y_i^{-1})\) of \(\mathbf{G} \times \mathbf{G}\) into \(\mathbf{G}\) is the composition of the mapping \(((x_i, y_i)) \to (x_i y_i^{-1})\) of

\[
\prod_ {i \in I} G _ {i} \times G _ {i} \text {   into   } G,
\]

and the canonical mapping \(((x_{i}), (y_{i})) \to ((x_{i}, y_{i}))\) of

\[
\mathbf {G} \times \mathbf {G} \quad \text { onto } \quad \prod_ {i \in \mathbf {I}} (\mathbf {G} _ {i} \times \mathbf {G} _ {i});
\]

and both these mappings are continuous (Chapter I, § 4, no. 1, Corollary 1 of Proposition 1, and Proposition 2).

DEFINITION 2. The topological group obtained by giving the product set

\[
\mathbf {G} = \prod_ {i \in \mathbf {I}} \mathbf {G} _ {i}
\]

the group structure which is the product of the group structures of the \(G_{i}\) and the topology which is the product of the topologies of the \(G_{i}\) is called the product of the topological groups \(G_{i}\) .

If \((\mathbf{J}_{\kappa})_{\kappa \in \mathbb{K}}\) is a partition of I, then G is isomorphic to the product of the topological groups \(\prod_{i \in J_{\kappa}} G_{i}\) (associativity of the product).

If \(\mathbf{H}_{i}\) is a subgroup of \(\mathbf{G}_{i}\) , then the product of the topological groups \(\mathbf{H}_{i}\) is isomorphic to the subgroup \(\prod_{i} \mathbf{H}_{i}\) of \(\prod_{i} \mathbf{G}_{i}\) . In particular, if \(J\) is any subset of \(I\) , and \(J' = C_{J}\) , then the topological group \(\prod_{i \in J} \mathbf{G}_{i}\) is isomorphic to the normal subgroup \(\mathbf{G}_{J}' = (\prod_{i \in J} \mathbf{G}_{i}) \times (\prod_{i' \notin J} \{e_{i'}\})\) of \(\mathbf{G}\) . Since the projection of every open set is an open set, the projection \(\mathbf{pr}_{J}\) of \(\mathbf{G}\) onto \(\prod_{i \in J} \mathbf{G}_{i}\) is a strict morphism, and consequently the quotient group \(\mathbf{G} / \mathbf{G}_{J}'\) is isomorphic to \(\mathbf{G}_{J'}'\) ; \(\mathbf{G}\) is isomorphic to the product \(\mathbf{G}_{J}' \times (\mathbf{G} / \mathbf{G}_{J}')\) .

PROPOSITION 25. Let \((\mathbf{G}_{i})_{i\in I}\) be a family of topological groups, and let H be the normal subgroup of \(\mathbf{G} = \prod_{i\in I}\mathbf{G}_i\) consisting of all \(x = (x_{i})\) such that the \(x_{i}\) are equal to the identity element \(e_i\) of \(\mathbf{G}_i\) except for a finite number of indices. Then the subgroup H is dense in G.

This is a particular case of Chapter I, § 4, no. 3, Proposition 8.

Let \((\mathbf{X}_i)_{i \in I}\) be a family of topological spaces, and for each \(i \in I\) let \(\mathbf{G}_i\) be a topological group operating continuously on \(\mathbf{X}_i\) . It is clear that the product group \(\mathbf{G} = \prod_{i \in I} \mathbf{G}_i\) then operates continuously on the product space \(\mathbf{X} = \prod_{i \in I} \mathbf{X}_i\) according to the law

\[
\left(\left(s _ {i}\right), \left(x _ {i}\right)\right)\rightarrow \left(s _ {i}. x _ {i}\right)
\]

(Chapter I, § 4, no. 1, Proposition 1, Corollary 1, and Proposition 2). Moreover the orbit under G of a point \(x = (x_{i})\) of X is the product of the orbits of the \(x_{i}\) (with respect to the groups \(G_{i}\) ). Let \(\varphi_{i}\) be the canonical mapping of \(X_{i}\) onto \(X_{i}/G_{i}\) , and let \(\varphi = (\varphi_{i})\) be the product mapping of X onto \(\prod_{i\in I}(X_{i}/G_{i})\) ; then the preceding remark shows that the bijection canonically associated with \(\varphi\) maps the orbit space X/G onto \(\prod_{i\in I}(X_{i}/G_{i})\) . Moreover:

PROPOSITION 26. The bijection of X/G onto \(\prod_{i\in I}(\mathbf{X}_i / \mathbf{G}_i)\) canonically associated with \((\varphi_i)\) is a homeomorphism.

For since the \(\varphi_{i}\) are surjective and open, it follows that \(\varphi = (\varphi_{i})\) is an open mapping (Chapter I, § 5, no. 3, Corollary to Proposition 8).

COROLLARY. Let \((\mathbf{G}_{i})_{i\in I}\) be a family of topological groups, and for each \(i\in I\) let \(\mathbf{H}_i\) be a normal subgroup of \(\mathbf{G}_i\) ; let \(\varphi_i\) denote the canonical mapping of \(\mathbf{G}_i\) onto \(\mathbf{G}_i / \mathbf{H}_i\) . Let \(\mathbf{G} = \prod_{i\in I}\mathbf{G}_i\) , \(\mathbf{H} = \prod_{i\in I}\mathbf{H}_i\) . Then the bijective homomorphism of \(\mathbf{G} / \mathbf{H}\) onto \(\prod_{i\in I}(\mathbf{G}_i / \mathbf{H}_i)\) associated with the continuous homomorphism \((x_i)\to (\varphi_i(x_i))\) is an isomorphism of topological groups.

For this homomorphism is an isomorphism of group structures.

Remark. If G is a commutative topological group, written additively, then the mapping \((x,y)\to x+y\) of \(G\times G\) onto G is a strict morphism. For since \((x+x')+(y+y')=(x+y)+(x'+y')\) it is a homomorphism of \(G\times G\) onto G; also it is continuous, and the image of a neighbourhood \(V\times V\) of the origin in \(G\times G\) under this mapping is the neighbourhood \(V+V\) of the origin in G.

\subsection{SEMI-DIRECT PRODUCTS}

Let L, N be two subgroups of a group G, such that LN = NL; then LN is a subgroup of G, since

\[
(\mathrm{LN}) (\mathrm{LN}) ^ {- 1} = \mathrm{LNN} ^ {- 1} \mathrm{L} ^ {- 1} = \mathrm{LNL} = \mathrm{LLN} = \mathrm{LN}
\]

Moreover, the mapping \(\varphi: (x, y) \to xy\) of \(N \times L\) into \(G\) is injective if and only if \(N \cap L = \{e\}\) . For it is clear that if \(\varphi\) is injective, then \(N \cap L = \{e\}\) ; and conversely, the relation \(x' y' = xy\) , where \(x, x' \in N\) and \(y, y' \in L\) , implies that \(x^{-1}x' = yy'^{-1} \in N \cap L\) ; hence if \(N \cap L = \{e\}\) , then \(\varphi\) is injective. Thus \(\varphi\) is bijective if and only if \(LN = G\) and \(N \cap L = \{e\}\) .

If N is a normal subgroup of G (or more generally is normal in some subgroup of G containing \(N \cup L\) ), the condition LN = NL is automatically satisfied. Moreover, for each \(y \in L\) the mapping \(\sigma_{y}: x \to yxy^{-1}\) is an automorphism of the group N, and for any two elements u, v of L we have

\[
\sigma_ {u v} = \sigma_ {u} \circ \sigma_ {v};\tag{1}
\]

and for any \(x, x'\) in \(\mathbf{N}\) and \(y, y'\) in \(\mathbf{L}\) we have also

\[
(x y) (x ^ {\prime} y ^ {\prime}) = (x \sigma_ {y} (x ^ {\prime})) (y y ^ {\prime}).\tag{2}
\]

Conversely :

PROPOSITION 27. Let N and L be two groups, \(e'\) and \(e''\) their respective identity elements. Suppose that we are given a homomorphism \(y \to \sigma_y\) of L into the group \(\Gamma\) of automorphisms of N. Then:

\begin{enumerate}
\def\labelenumi{\arabic{enumi})}
\tightlist
\item
  On the product set \(\mathbf{S} = \mathbf{N} \times \mathbf{L}\) , the internal law of composition
\end{enumerate}

\[
(x, y) (x ^ {\prime}, y ^ {\prime}) = (x \sigma_ {y} (x ^ {\prime}), y y ^ {\prime})\tag{3}
\]

defines a group structure, for which \(j_1: x \to (x, e'')\) is an isomorphism of \(N\) onto a normal subgroup of \(S\) , \(j_2: y \to (e', y)\) is an isomorphism of \(L\) onto a subgroup of \(S\) , and \(\mathrm{pr}_2: S \to L\) is a surjective homomorphism whose kernel is \(j_1(N)\) and which is such that \(\mathrm{pr}_2 \circ j_2\) is the identity automorphism of \(L\) .

\begin{enumerate}
\def\labelenumi{\arabic{enumi})}
\setcounter{enumi}{1}
\tightlist
\item
  Let \(f: \mathbf{N} \to \mathbf{G}\) , \(g: \mathbf{L} \to \mathbf{G}\) be two homomorphisms into a group \(\mathbf{G}\) , such that
\end{enumerate}

\[
f (\sigma_ {y} (x)) = g (y) f (x) g (y ^ {- 1})\tag{4}
\]

for all \(x \in \mathbb{N}\) and all \(y \in \mathbb{L}\) . Then there is a unique homomorphism \(h: \mathbb{S} \to \mathbb{G}\) such that \(f = h \circ j_1\) , and \(g = h \circ j_2\) .

If \(x, x', x''\) are elements of N and \(y, y', y''\) are elements of L, then

and

\[
\begin{array}{r l} ((x, y) (x ^ {\prime}, y ^ {\prime})) (x ^ {\prime \prime}, y ^ {\prime \prime}) & = (x \sigma_ {y} (x ^ {\prime}), y y ^ {\prime}) (x ^ {\prime \prime}, y ^ {\prime \prime}) \\ & = (x \sigma_ {y} (x ^ {\prime}) \sigma_ {y y ^ {\prime}} (x ^ {\prime \prime}), y y ^ {\prime} y ^ {\prime \prime}) \\ (x, y) ((x ^ {\prime}, y ^ {\prime}) (x ^ {\prime \prime}, y ^ {\prime \prime})) & = (x, y) (x ^ {\prime} \sigma_ {y ^ {\prime}} (x ^ {\prime \prime}), y ^ {\prime} y ^ {\prime \prime}) \\ & = (x \sigma_ {y} (x ^ {\prime} \sigma_ {y ^ {\prime}} (x ^ {\prime \prime})), y y ^ {\prime} y ^ {\prime \prime}) \end{array}
\]

and therefore the associativity of the law (3) follows from the facts that \(y \to \sigma_{y}\) is a homomorphism of L into \(\Gamma\) and that \(\sigma_{y}\) is an automorphism of N. Clearly \((e', e'')\) is the identity element of (3), and finally

\[
(x, y) \left(\sigma_ {y ^ {- 1}} (x ^ {- 1}), y ^ {- 1}\right) = \left(\sigma_ {y ^ {- 1}} (x ^ {- 1}), y ^ {- 1}\right) (x, y) = \left(e ^ {\prime}, e ^ {\prime \prime}\right)
\]

so that \((x, y)\) has an inverse in S. The other assertions of 1) are clear. On the other hand, since \((x, y) = (x, e'') (e', y)\) , a homomorphism \(h\) which satisfies the conditions of 2) must necessarily satisfy

\[
h (x, y) = f (x) g (y),
\]

hence is unique if it exists; moreover it is immediate from (4) that

\[
f \left(x \sigma_ {y} (x ^ {\prime})\right) g (y y ^ {\prime}) = f (x) g (y) f (x ^ {\prime}) g (y ^ {- 1}) g (y) g (y ^ {\prime}) = f (x) g (y) f (x ^ {\prime}) g (y ^ {\prime})
\]

which shows that \((x, y) \to f(x)g(y)\) is indeed a homomorphism of S into G satisfying the conditions of 2).

COROLLARY. The homomorphism \(h\) defined in 2) of Proposition 27 is injective if and only if \(f\) and \(g\) are injective and \(f(\mathbf{N}) \cap g(\mathbf{L}) = \{e\}\) ; and \(h\) is surjective if and only if \(f(\mathbf{N})g(\mathbf{L}) = \mathbf{G}\) .

Since \(h(x, y) = f(x)g(y)\) , the second assertion is obvious; moreover it follows from (4) that \(f(\mathbf{N})g(\mathbf{L}) = g(\mathbf{L})f(\mathbf{N})\) , and the first assertion follows from the remarks at the beginning of this sub-section.

The group S defined in Proposition 27 is said to be the external semi-direct product of N and L (relative to \(\sigma\) ); we shall generally identify N (resp. L) with the normal subgroup \(j_{1}(N)\) {[}resp. the subgroup \(j_{2}(L)\) {]} of S. If \(\sigma_{y}\) is the identity element of \(\Gamma\) for all \(y \in L\) , we retrieve the usual notion of the product of two groups.

Now let G be a group, and let L and N be two subgroups of G such that LN = NL and such that N is normal in NL, so that for each \(y \in L\) , \(\sigma_{y} : x \to yxy^{-1}\) is an automorphism of N, and \(y \to \sigma_{y}\) is a homomorphism of L into the automorphism group \(\Gamma\) of N. It follows then from Proposition 27 that if S is the external semi-direct product of N and L (relative to \(\sigma\) ), then \(h : (x, y) \to xy\) is a homomorphism of S into G. h is bijective if and only if we have \(N \cap L = \{e\}\) and NL = G (Corollary to Proposition 27); G is then said to be the semi-direct product of its normal subgroup N and its subgroup L, and we often identify G with S by means of h.

PROPOSITION 28. Let L, N be two topological groups, let \(y \to \sigma_y\) be a homomorphism of L into the group of automorphisms \(\Gamma\) of the (non-topological) group structure of N; and suppose that the mapping \((x, y) \to \sigma_y(x)\) of \(N \times L\) into N is continuous. Then:

\begin{enumerate}
\def\labelenumi{\arabic{enumi})}
\item
  On the external semi-direct product S of N and L, relative to \(\sigma\) , the product of the topologies of N and L is compatible with the group structure; the canonical injections \(j_{1}: \mathrm{N} \to \mathrm{S}\) and \(j_{2}: \mathrm{L} \to \mathrm{S}\) are isomorphisms of the topological groups N and L respectively onto the subgroups \(j_{1}(\mathrm{N})\) and \(j_{2}(\mathrm{L})\) of S, and \(pr_{2}\) is a strict morphism of S onto L.
\item
  Let \(f: \mathbf{N} \to \mathbf{G}\) and \(g: \mathbf{L} \to \mathbf{G}\) be two continuous homomorphisms into a topological group \(\mathbf{G}\) , satisfying (4); then the homomorphism
\end{enumerate}

\[
(x, y) \rightarrow f (x) g (y)
\]

of S into G is continuous.

This is an immediate consequence of the definitions and of the properties of the product topology.

The topological group S thus defined is said to be the external topological semi-direct product of N and L (relative to \(\sigma\) ); notice that the condition imposed on \(\sigma\) implies that L operates continuously on the left on N according to the external law \((x, y) \to \sigma_{y}(x)\) {[}no. 4{]}.

Now let G be a topological group and let N and L be two subgroups of G such that G is the semi-direct product of N and L, qua non-topological group; it is then clear that the mapping \((x, y) \to \sigma_{y}(x)\) is continuous on \(N \times L\) , and that the canonical bijective homomorphism

\[
h: (x, y) \rightarrow x y
\]

of S onto G is continuous. But this homomorphism is not necessarily bicontinuous; when it is bicontinuous, G is said to be the topological semi-direct product of N and L. For this to be so it is necessary and sufficient that, if \(p: G \to N\) and \(q: G \to L\) are the mappings which make correspond to \(z \in G\) the unique elements \(p(z) \in N\) and \(q(z) \in L\) such that \(z = p(z)q(z)\) , then one of the mappings p, q is continuous (in which case both are continuous). It comes to the same thing to say that the restriction to L of the canonical mapping \(G \to G/N\) is an isomorphism of the topological group L onto the topological group G/N.

\section{UNIFORM STRUCTURES ON GROUPS}

\subsection{THE RIGHT AND LEFT UNIFORMITIES ON A TOPOLOGICAL GROUP}

In a topological group G we can perceive the possibility of defining a notion of ``sufficiently near points'' and hence a uniform structure, by operating as follows: if x and y are any two points of G, we apply to both points the translation which sends one of them, say x, to the identity element e; the ``proximity'' of x and y is then evaluated in some sense by the neighbourhood V of e into which y is translated. This translation, which consists in multiplying both x and y by \(x^{-1}\) , can be carried out on the right or on the left, and we shall see that in either case we obtain effectively a uniformity on G which is compatible with the topology of G. Let us take the case in which the translations are performed on the right; then to each neighbourhood V of e there corresponds the set \(V_{d}\) of pairs \((x, y) \in \mathbf{G} \times \mathbf{G}\) such that \(yx^{-1} \in V\) . Let \(G_{d}\) be the family of sets \(V_{d}\) , as V runs through the neighbourhood filter B of e. Then \(G_{d}\) is a fundamental system of entourages (Chapter II, § I, no. I). For since \(e \in V\) , the diagonal \(\Delta\) of \(G \times G\) is contained in \(V_{d}\) for each \(V \in B\) , hence \(G_{d}\) is a filter base and satisfies axiom \((\mathrm{U}_{\mathrm{I}}^{\prime})\) ; since the relations \(yx^{-1} \in V\) and \(xy^{-1} \in V^{-1}\) are equivalent, we have \(\mathbf{V}_{d}^{-1} = (\mathbf{V}^{-1})_{d}\) .

hence \(\mathbf{V}_{d}^{-1}\in \mathfrak{G}_d\) by \((\mathrm{GV}_{\mathrm{II}})\) , so that \((\mathrm{U}_{\mathrm{II}}^{\prime})\) is satisfied; and finally, the relations \(zx^{-1}\in \mathrm{V}\) and \(yz^{-1}\in \mathrm{V}\) imply \(yx^{-1}\in \mathrm{V.V}\) ; hence \(\mathbf{V}_d\circ \mathbf{V}_d\) is contained in \((\mathrm{V.V})_d\) , and \((\mathrm{GV}_{\mathrm{I}})\) shows that \(\mathfrak{G}_d\) satisfies \((\mathrm{U}_{\mathrm{III}}^{\prime})\) .

The uniformity defined by \(\mathfrak{G}_d\) is compatible with the topology of \(\mathbf{G}\) , for the relations \(y \in \mathrm{V}_d(x)\) and \(y \in \mathrm{V}.x\) are equivalent by definition; in other words \(\mathrm{V}_d(x) = \mathrm{V}.x\) .

The argument is analogous when the translations are on the left, and we may therefore make the following definition:

DEFINITION 1. The right (resp. left) uniformity on a topological group G is the uniformity for which a fundamental system of entourages is obtained by making correspond to each neighbourhood V of the identity element e, the set \(V_{d}\) (resp. \(V_{s}\) ) of pairs \((x, y)\) such that \(y x^{-1} \in V\) (resp. \(x^{-1} y \in V\) ).

If V runs through a fundamental system of neighbourhoods of e, the sets \(V_{d}\) (resp. \(V_{s}\) ) form a fundamental system of entourages of the right (resp. left) uniformity.

To each proposition on the topology of a uniform space there corresponds a proposition on the topology of a group; the translation is made according to Definition 1 and the formulae \(\mathrm{V}_{d}(x)=\mathrm{V}.x,\quad\mathrm{V}_{d}(\mathrm{A})=\mathrm{V}.\mathrm{A},\quad\mathrm{V}_{s}(x)=x.\mathrm{V},\quad\mathrm{V}_{s}(\mathrm{A})=\mathrm{A}.\mathrm{V}\) , which are immediate consequences of the definition. For example, if A is any non-empty subset of G then we have (Chapter II, § 1, no. 2, Corollary 1 to Proposition 2)

\[
\overline {{{\mathrm{A}}}} = \bigcap_ {\mathbf {V} \in \mathfrak {B}} \mathrm{V}. \mathrm{A} = \bigcap_ {\mathbf {V} \in \mathfrak {B}} \mathrm{A}. \mathrm{V}.\tag{1}
\]

Again (Chapter II, § 1, no. 2, Corollary 3 to Proposition 2), every Hausdorff group is regular.

The right and left uniformities on a topological group are in general distinct (see Exercise 4). Obviously they coincide if the group is commutative, for then \(\mathbf{V}_d = \mathbf{V}_s\) ; also they coincide if the group is compact (Chapter II, § 4, no. 1, Theorem 1).

In general, we shall denote by \(G_{s}\) (resp. \(G_{d}\) ) the uniform space obtained by giving the set G its left (resp. right) uniformity.

PROPOSITION 1. The left and right translations are isomorphisms of the right uniformity onto itself.

As to the right translations, the result is clear, since the relation \(yx^{-1} \in V\) is equivalent to \((ya)(xa)^{-1} \in V\) {[}in other words, the mapping \((x, y) \to (xa, ya)\) leaves \(V_d\) fixed{]}. For the left translations, the result follows from \((\mathrm{GV}_{\mathrm{III}})\) ; for \(yx^{-1} \in V\) if and only if \((ay)(ax)^{-1} \in aVa^{-1}\) ; hence \(x \to ax\) is uniformly continuous on \(G_{d}\) .

Similarly, the right and left translations are isomorphisms of the left uniformity onto itself.

Every inner automorphism \(x \rightarrow axa^{-1}\) of G is therefore an automorphism for the group structure of G, for the topology of G, and for both the uniformities of G.

PROPOSITION 2. The symmetry \(x \to x^{-1}\) is an isomorphism of the right uniformity onto the left uniformity.

This is an immediate consequence of Definition 1.

\pandocbounded{\includegraphics[keepaspectratio]{images/part002_77aeac5d.jpg}}

The reader should beware of supposing that the mapping \((x, y) \to xy\) of the uniform space \(G_{d} \times G_{d}\) into the uniform space \(G_{d}\) is in general uniformly continuous. Similarly, the symmetry \(x \to x^{-1}\) , considered as a mapping of \(G_{d}\) onto \(G_{d}\) , is not in general uniformly continuous (see Exercises 3 and 4).

PROPOSITION 3. Every continuous homomorphism \(f\) of a topological group \(G\) into a topological group \(G'\) is uniformly continuous when considered as a mapping of \(G_d\) into \(G_d'\) (or of \(G_s\) into \(G_s'\) ).

For if \(V'\) is a neighbourhood of the identity element in \(G'\) , and \(\mathbf{V} = f^{-1}(\mathbf{V}')\) , then the relation \(yx^{-1} \in V\) implies

\[
f (y) (f (x)) ^ {- 1} = f (y x ^ {- 1}) \in V ^ {\prime}.
\]

\subsection{UNIFORMITIES ON SUBGROUPS, QUOTIENT GROUPS AND PRODUCT GROUPS}

If H is a subgroup of a topological group G, then the uniformity induced on H by the right uniformity of G is none other than the right uniformity of the topological group H.

If H is a normal subgroup of G, and if \(\varphi\) is the canonical mapping of G onto G/H, we obtain a fundamental system of entourages of the right uniformity of the quotient group G/H by associating with each neighbourhood V of the identity element in G, the set of all pairs \((\dot{x},\dot{y})\) of G/H such that \(\dot{x}\dot{y}^{-1}\in \varphi (\mathrm{V})\) (§ 2, no. 6, Proposition 17). This condition means that there is at least one point \(x\in \dot{x}\) and at least one point \(y\in \dot{y}\) such that \(yx^{-1}\in V\) {[}i.e.~such that \((x,y)\in V_d]\) . In particular, if N is the closure of the subset \(\{e\}\) of G, then the right uniformity on G/N is isomorphic to the Hausdorff uniformity associated with the right uniformity on G (cf.~Chapter II, § 3, no. 8).

Finally, on a product of a family \((\mathbf{G}_{i})\) of topological groups, the right uniformity is the product of the right uniformities of the \(G_{i}\) (cf.~Chapter II, § 2, no. 6).

There are analogous results for the left uniformity.

The left and right uniformities on the product group \(\prod_{i\in I}G_i\) are identical if and only if the left and right uniformities on each factor \(G_{i}\) coincide. This will always be the case if some of the \(G_{i}\) are commutative and the others are compact.

\subsection{COMPLETE GROUPS}

DEFINITION 2. A topological group is said to be complete if its left and right uniformities are structures of complete spaces.

From Proposition 2 of no. 1, for a group to be complete it is sufficient that one of its uniformities is a structure of a complete space. G is complete if and only if its associated Hausdorff group (§ 2, no. 6) is complete.

Every closed subgroup of a complete group is complete (Chapter II, § 3, no. 4, Proposition 8). Every product of complete groups is complete (Chapter II, § 3, no. 5, Proposition 10).

\pandocbounded{\includegraphics[keepaspectratio]{images/part002_8d186292.jpg}}

On the other hand, if G is a complete group and H is a closed normal subgroup of G, then the quotient group G/H is not necessarily complete (however, see Chapter IX, § 3, no. 1, Proposition 4).

PROPOSITION 4. If, in a topological group G, there is a neighbourhood V of e which is complete with respect to either the right or the left uniformity, then G is complete.

Suppose for example that V is complete with respect to the right uniformity, and let F be a Cauchy filter on \(G_{d}\) ; then F contains a \(V_{d}\) -small set M, and if \(x_{1} \in M\) we therefore have \(M \subset Vx_{1}\) . Hence the trace of F on the complete subspace \(Vx_{1}\) of \(G_{d}\) is a Cauchy filter, which converges to a point \(x_{0}\) ; since \(x_{0}\) is a cluster point of F, it is a limit of F (Chapter II, § 3, no. 2, Corollary 2 to Proposition 5).

COROLLARY 1. A locally compact group is complete.

For every compact space is complete with respect to its unique uniformity (Chapter II, § 4, no. 1, Theorem 1).

COROLLARY 2. Every locally compact subgroup of a Hausdorff topological group G is closed in G.

For every complete subspace of a Hausdorff uniform space is closed (Chapter II, § 3, no. 4, Proposition 8).

PROPOSITION 5. Let \(G_{1}\) be a topological group, let \(G_{2}\) be a complete Hausdorff topological group, and let \(H_{1}\) (resp. \(H_{2}\) ) be a dense subgroup of \(G_{1}\) (resp. \(G_{2}\) ). Then every continuous homomorphism u of \(H_{1}\) into \(H_{2}\) can be uniquely extended to a continuous homomorphism \(\overline{u}\) of \(G_{1}\) into \(G_{2}\) . Furthermore, if \(G_{1}\) is Hausdorff and complete, and if u is an isomorphism of \(H_{1}\) onto \(H_{2}\) , then \(\overline{u}\) is an isomorphism of \(G_{1}\) onto \(G_{2}\) .

\(\overline{u}\) is uniformly continuous with respect to the right uniformities of \(\mathbf{H}_1\) and \(\mathbf{H}_2\) (no. 1, Proposition 3), hence admits a unique extension to a mapping \(\overline{u}\) of \(\mathbf{G}_1\) into \(\mathbf{G}_2\) which is uniformly continuous with respect to the right uniformities of these groups (Chapter II, § 3, no. 6, Theorem 2). Moreover, by virtue of the principle of extension of identities (Chapter I, § 8, no. 1, Corollary 1 to Proposition 2), \(\overline{u}\) is a homomorphism of \(\mathbf{G}_1\) into \(\mathbf{G}_2\) , whence the first assertion of the Proposition. To prove the second assertion, it is enough to consider the isomorphism \(v\) of \(\mathbf{H}_2\) onto \(\mathbf{H}_1\) , which is the inverse of \(u\) , and its extension \(\overline{v}\) to a continuous homomorphism of \(\mathbf{G}_2\) into \(\mathbf{G}_1\) ; by reason of the uniqueness of the extension, \(\overline{v} \circ \overline{u}\) and \(\overline{u} \circ \overline{v}\) are the identity mappings of \(\mathbf{G}_1\) and \(\mathbf{G}_2\) respectively, and therefore (Set Theory, R, § 2, no. 12) \(\overline{u}\) is bijective.

\pandocbounded{\includegraphics[keepaspectratio]{images/part002_2e6b534c.jpg}}

Remark. If the continuous homomorphism \(u\) is bijective, it does not follow in general that \(\overline{u}\) is either injective or surjective (cf.~Exercise 12); but see no. 5, Proposition 9.

\subsection{COMPLETION OF A TOPOLOGICAL GROUP}

Let G be a Hausdorff topological group. The uniform space \(G_{d}\) may be considered as a dense subspace of its completion \(\hat{G}_{d}\) . We shall investigate whether G can be considered as a dense subgroup of a complete Hausdorff group \(G'\) . If so, then the uniform space \(G_{d}'\) must be isomorphic to \(\hat{G}_{d}\) (Chapter II, § 3, no. 6, Corollary to Theorem 2), and we must therefore be able to define on \(\hat{G}_{d}\) a topological group structure which induces the given topological group structure on G. Consequently we have to examine: 1) Whether we can extend by continuity the functions xy and \(x^{-1}\) to \(\hat{G}_{d} \times \hat{G}_{d}\) and \(\hat{G}_{d}\) respectively; 2) whether the functions thus extended do indeed define a group structure on \(\hat{G}_{d}\) (they will then necessarily define a topological group structure on \(\hat{G}_{d}\) inducing the given structure on G). We have next to establish that 3) when the preceding operations are possible, the topological group which they define is complete. Finally, we shall see that 4) if there is a complete group satisfying the given conditions, then it is unique up to isomorphism.

\begin{enumerate}
\def\labelenumi{\arabic{enumi})}
\tightlist
\item
  Extension of \(xy\) and \(x^{-1}\) by continuity. Since the functions \(xy\) and \(x^{-1}\) are not in general uniformly continuous, we cannot apply the theorem of extension of uniformly continuous functions (Chapter II, § 3, no. 6, Theorem 2). Nevertheless, we can extend \(xy\) , by virtue of Proposition 11 of Chapter II, § 3, no. 6 and the following Proposition:
\end{enumerate}

PROPOSITION 6. Let \(\mathfrak{F}\) and \(\mathfrak{G}\) be two Cauchy filters on \(\mathbf{G}_d\) . Then the image of the filter \(\mathfrak{F} \times \mathfrak{G}\) under the mapping \((x, y) \to xy\) is a Cauchy filter base on \(\mathbf{G}_d\) . Let us evaluate the ``proximity'' of \(xy\) and \(x'y'\) in \(\mathbf{G}_d\) ; in other words, let us form the product \((x'y')(xy)^{-1} = x'y'y^{-1}x^{-1}\) . For each \(a \in \mathbf{G}\) , we can also write \((x'y')(xy)^{-1} = (x'a^{-1})(ay'y^{-1}a^{-1})(ax^{-1})\) . We shall see that by suitable choice of \(a\) , each of the three factors of this product is very small whenever the pairs \((x, y)\) and \((x', y')\) belong to a sufficiently small set of \(\mathfrak{F} \times \mathfrak{G}\) . Let V be any neighbourhood of \(e\) in G; then there is a \(\mathbf{V}_d\) -small set A ∈ \(\mathfrak{F}\) . Choose \(a\) in A, then if \(x\) and \(x'\) are any two points of A, we have \(x'a^{-1} \in V\) and \(ax^{-1} \in V\) . On the other hand, the relation \(ay'y^{-1}a^{-1} \in V\) is equivalent to

\[
y ^ {\prime} y ^ {- 1} \in a ^ {- 1} \mathrm{V}a = \mathrm{W},
\]

and since W is a neighbourhood of e, there is a \(W_{d}\) -small set \(B \in \mathfrak{G}\) . Hence for all \((x, y)\) and \((x', y')\) in A × B, we have \((x'y')(xy)^{-1} \in V^{3}\) , and this completes the proof.

In order that \(x^{-1}\) can be extended by continuity to \(\hat{G}_d\) it is necessary and sufficient that the image, under the symmetry \(x \to x^{-1}\) , of a Cauchy filter on \(G_d\) is a Cauchy filter on \(G_d\) (Chapter II, § 3, no. 6, Proposition 11). There are examples of topological groups in which this condition is not satisfied (cf.~Chapter X, § 3, Exercise 16); we shall suppose that it is satisfied for the remainder of this proof.

\begin{enumerate}
\def\labelenumi{\arabic{enumi})}
\setcounter{enumi}{1}
\item
  The extended functions \(xy\) and \(x^{-1}\) define a group structure on \(\hat{G}_d\) . For if we apply the principle of extension of identities (Chapter I, § 8, no. 1, Proposition 2, Corollary 1) to the functions \(x(yz)\) and \((xy)z\) , defined on \(\hat{G}_d \times \hat{G}_d \times \hat{G}_d\) and equal on the dense subspace \(G_d \times G_d \times G_d\) , we see that the law of composition \((x, y) \to xy\) is associative on \(\hat{G}_d\) . For the same reason, the functions \(x, ex, xe\) are identical on \(\hat{G}_d\) , and the functions \(e, xx^{-1}, x^{-1}x\) are identical on \(\hat{G}_d\) .
\item
  The topological group \(\hat{\mathbf{G}}_d\) is complete. Let \(\mathfrak{U}_d\) be its right uniformity, and let \(\mathfrak{U}\) be the uniformity on \(\hat{\mathbf{G}}_d\) obtained by completing the right uniformity of \(\mathbf{G}\) . Then \(\mathfrak{U}\) and \(\mathfrak{U}_d\) induce the same uniformity on \(\mathbf{G}\) , and therefore every Cauchy filter base \(\mathfrak{B}\) on \(\mathbf{G}\) with respect to \(\mathfrak{U}_d\) is also a Cauchy filter base with respect to \(\mathfrak{U}\) . Now \(\mathfrak{B}\) converges in \(\hat{\mathbf{G}}_d\) , since \(\mathcal{U}\) is a complete uniformity; and since \(\mathcal{U}\) and \(\mathcal{U}_d\) induce the same topology on \(\hat{\mathbf{G}}_d\) , it follows (Chapter II, § 3, no. 4, Proposition 9) that \(\mathcal{U}_d\) is a complete uniformity. This conclusion also shows that \(\mathcal{U}\) and \(\mathcal{U}_d\) coincide (Chapter II, § 3, no. 6, Corollary to Theorem 2).
\item
  Uniqueness. This follows from Proposition 5 of no. 3.
\end{enumerate}

To sum up, we have proved the following theorem:

THEOREM 1. A Hausdorff topological group G is isomorphic to a dense subgroup of a complete group \(\hat{G}\) if and only if the image, under the symmetry \(x \to x^{-1}\) , of a Cauchy filter with respect to the right uniformity of G is a Cauchy filter with respect to this uniformity. The complete group \(\hat{G}\) (which is called the completion of G) is then unique (up to isomorphism).

PROPOSITION 7. Let G be a Hausdorff topological group which has a completion \(\hat{G}\) . Then the closures in \(\hat{G}\) of the neighbourhoods of the identity element in G form a fundamental system of neighbourhoods of the identity element in \(\hat{G}\) .

Since \(\hat{G}\) is regular, every neighbourhood of the identity element in \(\hat{G}\) contains the closure V of an open neighbourhood U of e in \(\hat{G}\) , and V is also the closure of the trace of U on G.

Let G be a group which is not necessarily Hausdorff; let \(N = \overline{e}\) , and let \(G' = G/N\) be the Hausdorff group associated with G (§ 2, no. 6). If \(G'\) has a completion \(\hat{G}'\) , this completion is called the Hausdorff completion of G and is denoted by \(\hat{G}\) ; \(\hat{G}_{d}'\) (resp. \(\hat{G}_{s}'\) ) is then the Hausdorff completion (Chapter II, § 3, no. 7) of the uniform space \(G_{d}\) (resp. \(G_{s}\) ).

PROPOSITION 8. Let G be a topological group which has a Hausdorff completion \(\hat{G}'\) . Then every continuous homomorphism \(u\) of G into a complete Hausdorff group H can be uniquely factorized into \(u = v \circ \varphi\) , where \(v\) is a continuous homomorphism of \(\hat{G}'\) into H and \(\varphi\) is the canonical mapping of G into \(\hat{G}'\) (the composition of the canonical injection of G' into \(\hat{G}'\) and the canonical homomorphism \(\psi\) of G onto G/N = G').

Since the kernel of \(u\) is closed and contains \(e\) , it contains \(N\) , and hence \(u\) can be written as \(u = w \circ \psi\) , where \(w\) is a continuous homomorphism of \(G'\) into \(H\) ; now apply Proposition 5 of no. 3 to \(w\) .

\subsection{UNIFORMITY AND COMPLETION OF A COMMUTATIVE TOPOLOGICAL GROUP}

We have already remarked that the left and right uniformities coincide on a commutative topological group G; whenever we speak of the uniformity of G, it is this unique uniformity to which we refer.

THEOREM 2. Let G be a commutative topological group. Then the functions \(x^{-1}\) and xy are uniformly continuous on G and G × G respectively. Moreover G admits a Hausdorff completion \(\hat{G}\) , and \(\hat{G}\) is commutative.

The uniform continuity of \(x^{-1}\) follows from Proposition 2 of no. 1, and that of xy from Proposition 3 of no. 1, since \((x, y) \to xy\) is a continuous homomorphism of \(G \times G\) into G. If G is Hausdorff, it satisfies the condition of Theorem 1 of no. 4 (as does every Hausdorff group whose left and right uniformities coincide); moreover the functions xy and yx are equal on \(\hat{G} \times \hat{G}\) by the principle of extension of identities; hence the second part of the theorem, by considering in the general case the Hausdorff group associated with G.

It follows in particular from this theorem that if \(f\) and \(g\) are two uniformly continuous mappings of a uniform space \(\mathbf{X}\) into a commutative group \(G\) , written additively, then the functions \(-f\) and \(f + g\) are uniformly continuous.

PROPOSITION 9. Let G be a commutative group, and let \(\mathfrak{T}_1\) , \(\mathfrak{T}_2\) be two Hausdorff topologies compatible with the group structure of G. Suppose that \(\mathfrak{T}_1\) is finer than \(\mathfrak{T}_2\) and that there is a fundamental system of neighbourhoods of o for \(\mathfrak{T}_1\) which are closed for \(\mathfrak{T}_2\) . Let \(\mathrm{G}_1, \mathrm{G}_2\) be the completions of G with respect to the topologies \(\mathfrak{T}_1, \mathfrak{T}_2\) respectively, and let \(f: \mathrm{G}_1 \to \mathrm{G}_2\) be the continuous homomorphism which extends the identity mapping of G (no. 3, Proposition 5). Then \(f\) is injective.

Suppose that G is written additively. Let \(\mathfrak{U}_{1}\) be the uniformity on G corresponding (no. 1) to the topology \(\mathfrak{T}_{1}\) : it will suffice to show that if F and \(F'\) are two minimal Cauchy filters (Chapter II, § 3, no. 2) with respect to \(\mathfrak{U}_{1}\) , which converge in \(G_{2}\) to the same point a, then \(F = F'\) (Chapter II, § 3, no. 7). For this, it is enough to show that \(F \cap F'\) is a Cauchy filter with respect to \(\mathfrak{U}_{1}\) . Let V be a neighbourhood of o in G with respect to \(\mathfrak{T}_{1}\) , such that V is closed in \(\mathfrak{T}_{2}\) , and let W be a symmetric neighbourhood of o in G with respect to \(\mathfrak{T}_{1}\) , such that \(W + W \subset V\) . By hypothesis, there is a \(W_{d}\) -small set M (resp. \(M'\) ) in F (resp. \(F'\) ); if \(x \in M\) and \(y \in M\) we have \(y - x \in W\) , i.e.~\(y \in x + W\) . If \(\overline{W}\) and \(\overline{V}\) are the closures of W and V in \(G_{2}\) , it follows that \(y \in x + \overline{W}\) , and therefore, since a is in the closure of M, that \(a \in x + \overline{W}\) for each \(x \in M\) . Similarly, \(a \in x' + \overline{W}\) for each \(x' \in M'\) , and hence \(x - x' \in \overline{W} + \overline{W}\) ; but since \((x, y) \to x + y\) is a continuous mapping of \(G_{2} \times G_{2}\) into \(G_{2}\) , we have \(\overline{W} + \overline{W} \subset \overline{W + W} \subset \overline{V}\) . It follows that if \(x \in M\) and \(x' \in M'\) , then \(x - x' \in \overline{V} \cap G = V\) , because V is closed in \(\mathfrak{T}_{2}\) ; and this completes the proof.

COROLLARY 1. Under the hypotheses of Proposition 9, if A is a subset of G which is a complete subspace with respect to the uniformity \(\mathfrak{U}_{2}\) corresponding to \(\mathfrak{T}_{2}\) , then A is also a complete subspace with respect to the uniformity \(\mathfrak{U}_{1}\) corresponding to \(\mathfrak{T}_{1}\) .

If \(A_{1}\) is the closure of A in \(G_{1}\) , then \(f(A_{1})\) is contained in the closure of A in \(G_{2}\) , which by hypothesis is equal to A. Since \(f(A)=A\) by definition and f is injective, we have \(A_{1}=A\) .

COROLLARY 2. Let G be a commutative group and let \(\mathfrak{T}_1, \mathfrak{T}_2\) be two topologies compatible with the group structure of G. Suppose that \(\mathfrak{T}_1\) is finer than \(\mathfrak{T}_2\) and that there is a fundamental system \(\mathfrak{B}\) of neighbourhoods of o with respect to \(\mathfrak{T}_1\) which are complete with respect to the uniformity \(\mathfrak{U}_2\) corresponding to \(\mathfrak{T}_2\) . Then G is complete with respect to the uniformity \(\mathfrak{U}_1\) corresponding to \(\mathfrak{T}_1\) .

The sets of \(\mathfrak{B}\) are closed in the topology \(\mathfrak{T}_{2}\) , hence complete for the uniformity \(\mathfrak{U}_{1}\) by Corollary 1; the result therefore follows from Proposition 4 of no. 3.

\section{GROUPS OPERATING PROPERLY ON A TOPOLOGICAL SPACE; COMPACTNESS IN TOPOLOGICAL GROUPS AND SPACES WITH OPERATORS}

\subsection{GROUPS OPERATING PROPERLY ON A TOPOLOGICAL SPACE}

DEFINITION 1. Let G be a topological group operating continuously on a topological space X. G is said to operate properly on X if the mapping

\[
\theta : (s, x) \rightarrow (x, s. x) \text{ of } \mathbf{G} \times \mathbf{X} \text{ into } \mathbf{X} \times \mathbf{X}
\]

is proper (Chapter I, § 10, no. 1, Definition 1).

Let \(\Gamma \subset \mathbf{G} \times \mathbf{X} \times \mathbf{X}\) be the graph of the mapping \(\rho: (s, x) \to s.x\) . Since \(\rho\) is continuous, the mapping \(\sigma: (s, x) \to (s, x, s.x)\) is a homeomorphism of \(\mathbf{G} \times \mathbf{X}\) onto \(\Gamma\) , and the composite mapping

\[
\mathbf {G} \times \mathbf {X} \xrightarrow {\sigma} \Gamma \xrightarrow {\mathrm{pr} _ {2 3}} \mathbf {X} \times \mathbf {X}
\]

is just \(\theta\) . Definition 1 is therefore equivalent to saying that the restriction of \(\operatorname{pr}_{23}\) to \(\Gamma\) is a proper mapping of \(\Gamma\) into \(\mathbf{X} \times \mathbf{X}\) .

Theorem 1 of Chapter I, § 10, no. 2 shows that G operates properly on X if and only if the following condition is satisfied:

For each set A filtered by an ultrafilter F, and each mapping

\[
\alpha \rightarrow (s _ {\alpha}, x _ {\alpha})
\]

of A into \(\mathbf{G} \times \mathbf{X}\) , if the mapping \(\alpha \to (s_{\alpha}.x_{\alpha}, x_{\alpha})\) has a limit \((b, a)\) with respect to \(\mathfrak{F}\) , then \(\alpha \to s_{\alpha}\) has a limit \(t \in \mathbf{G}\) with respect to \(\mathfrak{F}\) , such that \(t.a = b\) .

Examples. 1) Let H be a closed subgroup of a topological group G. If G operates properly on X, then so does H, since H × X is closed in G × X (Chapter I, § 10, no. 1, Proposition 5, Corollary 1). If for example we take X = G, with G operating on itself by left translations, then the mapping G × X → X × X is a homeomorphism, hence proper; thus H operates properly on G by left translations.

\begin{enumerate}
\def\labelenumi{\arabic{enumi})}
\setcounter{enumi}{1}
\tightlist
\item
  If G operates properly on X, then it operates properly on every subspace \(X'\) of X which is a union of orbits of points of X (in other words, \(X'\) is saturated with respect to the equivalence relation defined by G). For the inverse image of \(X' \times X'\) in \(G \times X\) is \(G \times X'\) , and we can apply Proposition 3 of Chapter I, § 10, no. 1.
\end{enumerate}

PROPOSITION 1. Let G be a topological group operating continuously on a topological space X, and let K be a quasi-compact subset of G. Then the mapping \(\rho: (s, x) \to s.x\) of \(K \times X\) into \(X\) is proper.

\(\rho\) factorizes into \(\mathbf{K} \times \mathbf{X} \to \mathbf{K} \times \mathbf{X} \xrightarrow{\mathrm{pr}_2} \mathbf{X}\) , where \(\alpha(s, x) = (s, s.x)\) . \(\alpha\) is a homeomorphism, for \(\alpha^{-1}: (s, y) \to (s, s^{-1}.y)\) is continuous. Since \(\mathbf{K}\) is quasi-compact, \(\mathrm{pr}_2\) is proper (Chapter I, § 10, no. 2, Theorem 1, Corollary 5); hence \(\rho\) is proper (Chapter I, § 10, no. 1, Proposition 5).

COROLLARY 1. If A is a closed (resp. compact) subset of X, then K.A is closed in X (resp. compact if X is Hausdorff).

The assertion concerning closed sets follows from Proposition 1 and the fact that a proper mapping is closed (Chapter I, § 10, no. 1, Proposition 1). The assertion concerning compact sets is trivial.

\pandocbounded{\includegraphics[keepaspectratio]{images/part002_62417fce.jpg}}

It should be noted that if L is a compact subset of X, and F a closed subset of G, then F. L is not necessarily closed in X (§ 2, Exercise 29; cf.~no. 5, Corollary to Proposition 12).

COROLLARY 2. If K is a quasi-compact subgroup of a topological group G, then the equivalence relation \(x^{-1}y \in K\) is closed, and the canonical mapping \(\varphi: G \to G/K\) is proper.

The first assertion follows from Corollary 1 applied to G acting by right translations on itself, and the second assertion then follows from Chapter I, § 10, no. 2, Theorem 1.

COROLLARY 3. Let K be a quasi-compact normal subgroup of a topological group G, and let \(\varphi\) be the canonical mapping G → G/K. Then for each closed subgroup A of G the canonical bijection of A/A∩K onto \(\varphi(A)\) is an isomorphism of topological groups.

Since \(x^{-1}y \in K\) is a closed equivalence relation (Corollary 2), the Corollary follows from Chapter I, § 5, no. 2, Proposition 4.

PROPOSITION 2. Let K be a compact group operating continuously on a Hausdorff space X. Then:

\begin{enumerate}
\def\labelenumi{\alph{enumi})}
\item
  K operates properly on X.
\item
  The mapping \((s, x) \to s.x\) of \(\mathbf{K} \times \mathbf{X}\) into \(\mathbf{X}\) is proper.
\item
  The canonical mapping of X onto X/K is proper.
\item
  is a consequence of Proposition 1. As to a), since K is compact, \(\operatorname{pr}_{2}:(s,x)\to x\) is proper (Chapter I, § 10, no. 2, Theorem 1, Corollary 5); hence, X being Hausdorff, \((s,x)\to(x,s.x)\) is proper (Chapter I, § 10, no. 1, Proposition 5, Corollary 3). It remains to prove c). By Corollary 1 to Proposition 1, the canonical mapping \(\varphi:X\to X/K\) is closed. If Z is any topological space and if we make K operate trivially on Z, then K operates continuously on \(X\times Z\) and hence the canonical mapping \(X\times Z\to(X\times Z)/K\) is closed. But \((X\times Z)/K\) can be canonically identified with \((X/K)\times Z\) (§ 2, no. 4, Lemma 2 and Chapter I, § 5, no. 3, Corollary to Proposition 8); hence the canonical mapping \(X\times Z\to(X\times Z)/K\) can be identified with \(\varphi\times1\) , and since it is closed for all Z it follows that \(\varphi\) is proper.
\end{enumerate}

COROLLARY 1. Under the hypotheses of Proposition 2, X is compact (resp. locally compact) if and only if X/K is compact (resp. locally compact).

This follows from the fact that the canonical mapping \(X \rightarrow X/K\) is proper, in view of Proposition 9 of Chapter I, § 10, no. 4.

COROLLARY 2. Let G be a Hausdorff topological group and let K be a compact subgroup of G. Then G is compact (resp. locally compact) if and only if G/K is compact (resp. locally compact).

Apply Corollary 1 to K operating on G by right translations.
